\clearpage
\fancyhead[L]{\small VESTA. CERES. KOMET.}
\SourceLeaf{313}\begin{center}
VESTA. CERES. KOMET.
\end{center}

\par\vspace{4pt}\hrule\vspace{4pt}
\begin{center}
\fontsize{10}{12}\selectfont Göttingische gelehrte Anzeigen. Stück 180. Seite 1793..1795. 10. November 1808.
\end{center}
\par\vspace{4pt}\hrule\vspace{4pt}

Professor \textsc{Bode} hat uns die folgenden wertvollen Beobachtungen der \emph{Vesta} mitgeteilt, die an der Berliner Sternwarte gemacht wurden:


\begin{center}\fontsize{10}{12}\selectfont\setlength{\tabcolsep}{5pt}\renewcommand{\arraystretch}{1.15}\begin{tabular}{lrrr}
1808 & \shortstack{Mittlere Zeit\\in Berlin} & \shortstack{Scheinbare\\Rektaszension} & \shortstack{Scheinbare südliche\\Deklination}\\
\hline
Aug. 24 & $11^{\text{u}}\;16^{\text{m}}\;16^{\text{s}}$ & $354^\circ\;34'\;59''$ & $13^\circ\;50'\;22''$\\
Sept. 1 & $12\;48\;0.7$ & $352\;57\;27$ & $14\;49\;41$\\
5 & $11\;8\;0$ & $352\;5\;20$ & $15\;20\;11$\\
13 & $11\;50\;1.0$ & $350\;15\;10$ & $16\;15\;59$\\
15 & $11\;40\;21.7$ & $349\;48\;2$ & $16\;28\;2$\\
20 & $11\;16\;26.7$ & $348\;42\;5$ & $16\;54\;35$\\
Oct. 6 & $10\;1\;55.2$ & $345\;49\;31$ & $17\;42\;7$\\
9 & $9\;48\;36.0$ & $345\;26\;47$ & $17\;44\;4$\\
14 & $9\;26\;57.3$ & $344\;56\;58$ & $17\;42\;30$\\
\end{tabular}\end{center}


\par\indent Alle diese Beobachtungen bis auf die erste und die dritte wurden am Mauerquadranten gemacht.

\par\indent \emph{Ceres} blieb am Mauerquadranten stets unsichtbar. Die folgenden drei Beobachtungen wurden mit dem Kreismikrometer gemacht:


\begin{center}\fontsize{10}{12}\selectfont\setlength{\tabcolsep}{5pt}\renewcommand{\arraystretch}{1.15}\begin{tabular}{lrrr}
1808 & \shortstack{Mittlere Zeit\\in Berlin} & \shortstack{Scheinbare\\Rektaszension} & \shortstack{Scheinbare südliche\\Deklination}\\
\hline
Aug. 29 & $10^{\text{u}}\;52^{\text{m}}\;49^{\text{s}}$ & $314^\circ\;45'\;57''$ & $30^\circ\;34'\;28''$\\
Sept. 16 & $7\;46\;36$ & $312\;26\;31$ & $30\;39\;58$\\
17 & $9\;12\;43$ & $312\;22\;24$ & $30\;40\;4$\\
\end{tabular}\end{center}


\par\indent Die folgenden Beobachtungen des letzten Kometen wurden von den Herren \textsc{Schubert} und \textsc{Wisniewsky} in Petersburg angestellt und uns ebenfalls von Professor \textsc{Bode} freundlicherweise mitgeteilt. Sie sind von besonderem Interesse, weil es die spätesten Beobachtungen sind, die irgendwo gemacht werden konnten. Bei den letzten war der Komet so lichtschwach, dass ihn nur ein sehr scharfes Auge im Fernrohr erkennen konnte. Die Beobachtungen wurden mit dem Kreismikrometer gemacht. Bei der Reduktion berücksichtigte man Refraktion, Nutation und Aberration nicht. Die Vergleichssterne wurden nicht angegeben.


\begin{center}\fontsize{10}{12}\selectfont\setlength{\tabcolsep}{5pt}\renewcommand{\arraystretch}{1.15}\begin{tabular}{lrrr}
1808 & \shortstack{Mittlere Zeit\\in Petersburg} & \shortstack{Rektaszension} & \shortstack{Nordliche\\Deklination}\\
\hline
Jan. 15 & $6^{\text{u}}\;49^{\text{m}}\;30^{\text{s}}$ & $335^\circ\;50'\;2''$ & $47^\circ\;47'\;27''$\\
26 & $7\;2\;25$ & $345\;33\;24$ & $48\;4\;6$\\
Febr. 15 & $8\;45\;44$ & $2\;36\;58$ & $48\;44\;20$\\
25 & $9\;17\;3$ & $9\;21\;23$ & $48\;50\;48$\\
März 1 & $9\;53\;5$ & $12\;33\;34$ & $48\;51\;59$\\
22 & $9\;24\;29$ & $23\;6\;39$ & $48\;48\;25$\\
27 & $10\;54\;20$ & $25\;48\;51$ & $48\;53\;49$\\
\end{tabular}\end{center}


\par\indent Bei dieser Gelegenheit fügen wir auch die Beobachtungen der \emph{Vesta} bei, die in diesem Jahr an der hiesigen Sternwarte gemacht wurden. Die erste, die überhaupt die früheste in diesem Jahr ge\SourceLeaf{314}machte zu sein scheint, wurde bereits im 107. Stück dieser Zeitschrift [S. 303 dieses Bandes] mitgeteilt. Seit dem 10. August wurden die Beobachtungen am Mauerquadranten gemacht. Zufällige Hindernisse haben die Fortsetzung im Oktober unterbrochen.


\begin{center}\fontsize{10}{12}\selectfont\setlength{\tabcolsep}{5pt}\renewcommand{\arraystretch}{1.15}\begin{tabular}{lrrr}
1808 & \shortstack{Mittlere Zeit\\in Göttingen} & \shortstack{Scheinbare\\Rektaszension} & \shortstack{Scheinbare südliche\\Deklination}\\
\hline
Juni 22 & $13^{\text{u}}\;46^{\text{m}}\;25^{\text{s}}$ & $353^\circ\;12'\;17''$ & $9^\circ\;12'\;16''$\\
Juli  1 & $14\;38\;57$ & $354\;28\;3$ & $9\;9\;39$\\
2 & $12\;50\;2$ & $355\;7\;7$ & \\
13 & $13\;22\;37$ & $356\;35\;27$ & $9\;27\;40$\\
Aug. 10 & $14\;28\;53.7$ & $356\;36\;16.5$ & $11\;52\;50.3$\\
12 & $14\;20\;22.3$ & $356\;23\;22.3$ & $12\;8\;5.3$\\
14 & $14\;11\;23.4$ & $356\;9\;3.7$ & $12\;23\;25.7$\\
15 & $14\;6\;54.8$ & $356\;1\;6.0$ & $12\;31\;21.2$\\
22 & $13\;35\;4.7$ & $354\;56\;6.3$ & $13\;28\;20.7$\\
23 & $13\;30\;28.0$ & $354\;45\;24.1$ & $13\;36\;10.4$\\
29 & $13\;2\;11.5$ & $353\;35\;29.2$ & $14\;25\;25.8$\\
30 & $12\;57\;26.4$ & $353\;23\;12.9$ & $14\;33\;32.6$\\
Sept. 5 & $12\;28\;37.0$ & $352\;4\;33.3$ & $15\;20\;32.5$\\
8 & $12\;14\;6.6$ & $351\;23\;46.5$ & $15\;42\;22.3$\\
11 & $11\;59\;34.6$ & $350\;43\;7.3$ & $16\;3\;7.1$\\
12 & $11\;54\;42.9$ & $350\;28\;40.5$ & $16\;9\;34.9$\\
21 & $11\;11\;22.7$ & $348\;29\;10.5$ & $16\;59\;51.1$\\
\end{tabular}\end{center}


\Needspace{17\baselineskip}
\par\indent Herr \textsc{Tiarks} widmet sich bei uns mit Erfolg den mathematischen Wissenschaften und ist bereits sehr geübt im astronomischen Rechnen. Er hat einige dieser Beobachtungen mit den Elementen verglichen, die wir im 107. Stück dieser Anzeigen [S. 303 dieses Bandes] mitgeteilt haben. Daraus ergaben sich folgende Differenzen:


\begin{center}\fontsize{10}{12}\selectfont\setlength{\tabcolsep}{5pt}\renewcommand{\arraystretch}{1.15}\begin{tabular}{lrr}
 & \multicolumn{2}{c}{Differenz}\\
 & in der Länge & in der Breite\\\hline
Aug. 23 & $+4'\;7''4$ & $+43''0$\\
29 & $+4\;5.6$ & $+50.6$\\
30 & $+3\;53.9$ & $+45.0$\\
Sept. 5 & $+4\;4.5$ & $+45.2$\\
8 & $+4\;11.3$ & $+54.0$\\
11 & $+3\;31.6$ & $+38.3$\\
12 & $+4\;7.7$ & $+59.6$\\
21 & $+3\;37.0$ & $+55.1$\\
\end{tabular}\end{center}


